\chapter{Fever in Babies and Children: A Quick Guide}

\section*{Definition and Overview}
This chapter is a quick reference guide to fever in babies and children. Fever is a temperature of 38 degrees Celsius or higher, and in children it is most often caused by a viral infection that resolves on its own. The most important thing to assess is not the number on the thermometer but how the child looks and behaves: a child who is alert, drinking, and playing, even with a high temperature, is usually fine, while a child who is floppy, struggling to breathe, or very drowsy needs urgent care. This guide summarises the key steps for measuring, treating, and monitoring fever, and the danger signs that require immediate medical attention.

\section*{Epidemiology}
Fever is one of the most common reasons children see a doctor or attend an emergency department. Most children have several febrile illnesses in their first few years, mainly viral, and the great majority recover at home without medical treatment. Around 1 in 3 to 1 in 5 children see a health professional for fever each year. Serious bacterial infection causing fever is uncommon but occurs most often in babies under 3 months, whose immune systems are immature. Understanding normal patterns and knowing the warning signs helps parents manage fever confidently and seek help when it truly matters.

\section*{Aetiology and Pathophysiology}
\begin{itemize}
    \item \textbf{Viral infections:} the cause of most childhood fevers; colds, influenza, gastroenteritis, roseola, and hand, foot and mouth disease
    \item \textbf{Bacterial infections:} ear infections, tonsillitis, pneumonia, urinary tract infection, and rarely serious infections such as meningitis or sepsis
    \item \textbf{Immunisation:} a short-lived fever is a normal immune response after some vaccines
    \item \textbf{Teething:} can cause a mild temperature rise, but not usually a significant fever
    \item \textbf{Mechanism:} infection triggers immune cells to release signalling molecules that raise the brain's temperature set point; the fever helps fight the infection
    \item \textbf{The child's response:} the immune system of a young child is still learning, which is why fevers are frequent, brief, and usually harmless
    \item \textbf{Febrile seizures:} about 3--5\% of children aged 6 months to 5 years have a convulsion with a rapid temperature rise; these are frightening but almost always benign
\end{itemize}

\section*{Clinical Features}
\begin{itemize}
    \item A temperature of 38 degrees or higher, measured with a digital thermometer
    \item The child feels hot, may be flushed, and may shiver while the temperature is rising
    \item Reduced appetite, increased thirst, irritability, or clinginess
    \item A child who is otherwise alert, interested in play, drinking, and wetting nappies normally
    \item Symptoms pointing to the cause: cough, runny nose, ear tugging, vomiting, diarrhoea, or pain on passing urine
    \item Danger features: floppiness, difficulty waking, rapid or laboured breathing, a rash that does not fade under pressure, a bulging fontanelle in a baby, or a seizure
    \item In babies under 3 months: any fever is significant and needs medical assessment
\end{itemize}

\section*{Differential Diagnosis}
Fever in a child has many causes, and the assessment looks for the source.
\begin{itemize}
    \item Viral infections, which predominate and need no specific treatment
    \item Ear infection, tonsillitis, and chest infection, which may need antibiotics
    \item Urinary tract infection, an important cause in young children with fever without other symptoms
    \item Gastroenteritis with vomiting and diarrhoea
    \item Serious infections: meningitis, sepsis, and pneumonia, which must be recognised early
    \item Post-immunisation fever
    \item Heat-related illness in hot weather
    \item Non-infectious causes: inflammatory conditions and, rarely, malignancy
\end{itemize}

\section*{When to Seek Medical Care}
See a doctor if the child is under 3 months with a temperature of 38 degrees or more; if fever lasts more than 3 days; if the child is not drinking enough, is in pain, or is not their usual self; or if you are worried for any reason. Trust parental instinct: it is a valid reason to seek help.

\textbf{Call 000 (or local emergency services) for:}
\begin{itemize}
    \item A baby under 3 months with a fever
    \item A child who is floppy, unresponsive, or difficult to wake
    \item Difficulty breathing, fast breathing, or a blue tinge to the lips
    \item A rash that does not fade when pressed with a glass
    \item A seizure lasting more than 5 minutes, or any seizure in a baby under 6 months
    \item No wet nappies for 6--8 hours, sunken eyes, or a dry mouth
    \item Stiff neck, severe headache, or an unusual cry
    \item A child who is getting worse despite treatment
\end{itemize}

\section*{Diagnosis and Investigations}
\begin{itemize}
    \item \textbf{Measure accurately:} use a digital thermometer; for babies, under the arm or rectally is most reliable, and for older children, in the ear or mouth
    \item \textbf{History and examination:} the doctor checks for the source of the fever and assesses how unwell the child looks
    \item \textbf{Urine test:} recommended for young children, especially girls, with fever and no other obvious cause
    \item \textbf{Blood tests:} for children who look unwell, have prolonged fever, or are under 3 months
    \item \textbf{Chest X-ray:} when pneumonia is suspected
    \item \textbf{Lumbar puncture:} for suspected meningitis
    \item \textbf{Observation:} the child's response to treatment and their overall progress guide further care
\end{itemize}

\section*{Management}
\begin{itemize}
    \item \textbf{Comfort:} dress the child lightly, keep the room cool, offer fluids often, and continue breast or formula feeds
    \item \textbf{Paracetamol or ibuprofen:} for comfort, at the dose for the child's age and weight, at the correct intervals; use one medicine at a time
    \item \textbf{Never use aspirin} in children under 16
    \item \textbf{Do not over-treat:} a normal temperature is not the goal; a comfortable, drinking child is
    \item \textbf{Treat the cause:} antibiotics for bacterial infections; viral infections need time and fluids
    \item \textbf{Febrile seizures:} place the child on their side, protect the head, time the seizure, and call for help if it lasts more than 5 minutes
    \item \textbf{Review:} see a doctor again if the child worsens, fever persists beyond 3 days, or new symptoms appear
\end{itemize}

\section*{Complications}
\begin{itemize}
    \item Dehydration from poor intake, especially with vomiting or diarrhoea
    \item Febrile seizures in susceptible children
    \item Progression of unrecognised bacterial infection, including sepsis and meningitis
    \item Medicine errors: double dosing, overlapping paracetamol and ibuprofen, or adult doses given to children
    \item Unnecessary worry when the height of the fever is mistaken for severity of illness
\end{itemize}

\section*{Prognosis and Outcomes}
Almost all childhood fevers settle within 3 days as the underlying infection resolves, and the fever itself does not cause harm or brain damage. Children with serious bacterial infection recover fully with prompt antibiotics, and the risk is lowest in children who look well between fevers. Febrile seizures do not cause lasting harm and most children outgrow them. With clear guidance and early review when needed, the great majority of febrile children are managed safely at home, and emergency care is reserved for the small number with warning signs.

\section*{Prevention and Self-Care}
\begin{itemize}
    \item Keep a digital thermometer and correct-dose paracetamol or ibuprofen at home
    \item Know your child's normal behaviour so you can judge when they are truly unwell
    \item Give fluids frequently and keep the child cool and comfortable
    \item Record the temperature, fluids, nappy output, and the child's overall state
    \item Follow the immunisation schedule to prevent common causes of fever
    \item Wash hands and limit contact with unwell people during illness season
    \item Have the emergency numbers and a plan for urgent care, including your nearest emergency department
    \item If unsure, seek medical advice; it is always reasonable to have a fever assessed
\end{itemize}

\section*{Sources and Further Reading}
The following reputable sources provide further information for healthcare students and parents seeking reliable, current advice about fever in babies and children.
\begin{itemize}
    \item The Royal Children's Hospital Melbourne. \textit{Fever in children fact sheet.} https://www.rch.org.au/kidsinfo/
    \item Healthdirect Australia. \textit{Fever in children.} https://www.healthdirect.gov.au/fever-in-babies
    \item Raising Children Network. \textit{Fever.} https://raisingchildren.net.au/
    \item Sydney Children's Hospitals Network. \textit{Fever guidelines for parents.} https://www.schn.health.nsw.gov.au/
    \item NHS. \textit{High temperature (fever) in children.} https://www.nhs.uk/conditions/fever-in-children/
    \item World Health Organization. \textit{Childhood fever management.} https://www.who.int/
\end{itemize}
